\section{A two-shift reflection for generalized harmonic sums}
\label{mld:gh:sec:reflection}\label{mld:sec:harmonic}

The generalized harmonic order and the elementary-symmetric harmonic index
in the repository are different parameters. We use new notation to keep
them separate. Put
\[
 h_n^{(r)}(u)=\sum_{j=1}^n\frac1{(j-u)^r},\qquad
 Q_{p,r}=\sum_{n=1}^{\infty}\frac{(-1)^nH_n^{(r)}}{(2n+1)^p},
 \qquad p,r\geq1.
\]
Thus $h_n^{(r)}(0)=H_n^{(r)}$, and $Q_{p,1}=S_p$ in the manuscript.
In contrast, the manuscript's $S_{p,r}$ uses
$e_r(n)=[v^r]\prod_{j=1}^n(1+v/j)$; it is not $Q_{p,r}$.

The unshifted even-weight evaluations belong to the established theory
of linear Euler $R$-sums. In the normalization of Xu--Wang
\cite[Section~3.4, equation~(20)]{XuWang},
\begin{equation}\label{mld:gh:eq:normalization}
 R_{r,p}^{1,-1}=2^p Q_{p,r}.
\end{equation}
We derive a two-variable digamma resummation, its arbitrary mixed
derivatives, and a concrete rational-shift evaluation. The scalar parity
theorem is credited to that literature. The priority of the particular
two-variable formulation below has not been established.

\subsection{Ordinary convergence, including the first outer power}

Define
\begin{align}
 A(z,u)&=\sum_{n=1}^{\infty}
       \frac{(-1)^n h_n^{(1)}(u)}{2n+1-z},\label{mld:gh:eq:A}\\
 D(v)&=\frac14\left\{
       \psi\left(\frac{v+3}{4}\right)
       -\psi\left(\frac{v+1}{4}\right)\right\}.
       \label{mld:gh:eq:D}
\end{align}
Initially, $D(v)=\int_0^1x^v/(1+x^2)\,dx$ for $\Re v>-1$;
the displayed expression supplies its meromorphic continuation.
The first sum in \eqref{mld:gh:eq:A} is an ordinary convergent alternating
series.

\begin{lemma}[Local uniform convergence]\label{mld:gh:lem:convergence}
For positive integers $p,r$, the series
\[
 \sum_{n=1}^{\infty}
 \frac{(-1)^n h_n^{(r)}(u)}{(2n+1-z)^p}
\]
converges locally uniformly on
\[
 \Omega_+=
 \bigl(\mathbb C\setminus\{3,5,7,\ldots\}\bigr)
 \times\bigl(\mathbb C\setminus\{1,2,3,\ldots\}\bigr).
\]
It is holomorphic there and can be differentiated term by term any
finite number of times in $z$ and $u$. It is absolutely convergent
when $p\geq2$.
\end{lemma}
\begin{proof}
Fix a compact subset of $\Omega_+$. Uniformly on that subset,
$h_n^{(1)}(u)=O(1+\log n)$, while $h_n^{(r)}(u)=O(1)$ for $r\geq2$.
For
$f_n=h_n^{(r)}(u)(2n+1-z)^{-p}$, the exact difference is
\begin{align*}
 f_{n+1}-f_n
 ={}&\frac{1}{(n+1-u)^r(2n+3-z)^p}\\
 &+h_n^{(r)}(u)
       \left\{(2n+3-z)^{-p}-(2n+1-z)^{-p}\right\}.
\end{align*}
Consequently $f_n\to0$ uniformly and
$f_{n+1}-f_n=O(n^{-p-r}+(1+\log n)n^{-p-1})$ uniformly.
The sum of these differences is absolutely convergent, so grouping
successive even and odd indices proves local uniform convergence of
the original alternating series. Its individual terms are holomorphic.
The same estimates apply to every termwise derivative, because
derivatives increase $p$ or $r$. This proves the assertion, including
the conditionally convergent case $p=1$. The stated absolute convergence
follows directly from the bounds for $h_n^{(r)}$.
\end{proof}

\subsection{The exact bivariate identity}

\begin{theorem}[Two-shift reflection]\label{mld:gh:thm:reflection}
On
\[
 \Omega=
 \bigl(\mathbb C\setminus\{\pm3,\pm5,\pm7,\ldots\}\bigr)
 \times\bigl(\mathbb C\setminus(\mathbb Z\setminus\{0\})\bigr),
\]
one has
\begin{equation}\label{mld:gh:eq:reflection}
 \boxed{A(z,u)+A(-z,-u)=B(z,u),}
\end{equation}
where
\begin{align}
 B(z,u)={}&-\frac{\pi}{2\cos(\pi z/2)}
   \left\{\psi(1-u)-\psi\left(\frac{1+z}{2}-u\right)\right\}
       \notag\\
 &+\frac{\pi D(z-2u)}{\sin\pi u}-\frac{D(z)}{u}.
       \label{mld:gh:eq:B}
\end{align}
Apparent singularities of the complete expression on $\Omega$ are
removable. This convention applies in particular at $u=0$, at $z=\pm1$,
and where an individual digamma term in \eqref{mld:gh:eq:B} has a pole.
There is no assertion that its summands can be evaluated separately
at such points.

For $|z|<3$ and $|u|<1$, the ordinary Taylor expansion of $A$ is
\begin{equation}\label{mld:gh:eq:Agenerating}
 A(z,u)=\sum_{p,r\geq1} Q_{p,r}\,z^{p-1}u^{r-1}.
\end{equation}
Thus $B$ determines precisely the terms with $p+r$ even in this germ.
\end{theorem}

\begin{proof}
We first prove the identity near $(z,u)=(0,0)$. Set
\[
 G_r(z)=\int_0^1
      \frac{x^{-z}\operatorname{Li}_r(-x^2)}{1+x^2}\,dx.
\]
The harmonic generating function
$\operatorname{Li}_r(w)/(1-w)=\sum_{n\geq1}H_n^{(r)}w^n$
gives
$G_r(z)=\sum_{p\geq1}Q_{p,r}z^{p-1}$.
One can first insert a common Abel factor and then pass to its limit;
the preceding lemma identifies that limit with the ordinary sum.
Alternatively, the logarithmic moment representation is
\begin{equation}\label{mld:gh:eq:moment}
 Q_{p,r}=\frac1{(p-1)!}\int_0^1
   \frac{(-\log x)^{p-1}\operatorname{Li}_r(-x^2)}{1+x^2}\,dx.
\end{equation}
For $0\leq y\leq1$ and integer $r\geq1$,
$|\operatorname{Li}_r(-y)|\leq y$ by the alternating-series test.
It follows from \eqref{mld:gh:eq:moment} that
$|Q_{p,r}|\leq3^{-p}$. Hence the double power series in
\eqref{mld:gh:eq:Agenerating} converges normally for the stated radii.
Its coefficients also follow by the justified derivatives of
\eqref{mld:gh:eq:A} at the origin.

The full Mellin identity, after the substitution $y=x^2$, is
\begin{equation}\label{mld:gh:eq:fullMellin}
 M_r(z):=\int_0^\infty
     \frac{x^{-z}\operatorname{Li}_r(-x^2)}{1+x^2}\,dx
 =-\frac{\pi}{2\cos(\pi z/2)}
   \left\{\zeta\left(r,\frac{1+z}{2}\right)-\zeta(r)\right\}.
\end{equation}
It initially suffices to take $|\Re z|<1$. At $r=1$, the expression
in braces denotes the finite difference
$\psi(1)-\psi((1+z)/2)$.

For $L\in\mathbb R$, define the polynomial
\begin{equation}\label{mld:gh:eq:Pr}
 P_r(L)=\frac{(2L)^r}{r!}
     +2\sum_{j=1}^{\lfloor r/2\rfloor}
       \eta(2j)\frac{(2L)^{r-2j}}{(r-2j)!},
 \qquad \eta(s)=(1-2^{1-s})\zeta(s).
\end{equation}
Then
\begin{equation}\label{mld:gh:eq:inversion}
 \operatorname{Li}_r(-e^{2L})
   +(-1)^r\operatorname{Li}_r(-e^{-2L})=-P_r(L).
\end{equation}
Here the signs and branch convention can be checked without invoking
a complex inversion formula. The derivative of the left side is twice
the corresponding expression for $r-1$. At $r=0$ that expression
equals $-1$. At $L=0$ it is zero for odd $r$, and $-2\eta(r)$ for
even positive $r$. These data give exactly \eqref{mld:gh:eq:Pr}.

Splitting \eqref{mld:gh:eq:fullMellin} at $1$ and substituting $x\mapsto1/x$
in its second part gives
\begin{equation}\label{mld:gh:eq:integerreflection}
 G_r(z)-(-1)^rG_r(-z)
 =M_r(z)+(-1)^r\int_0^1\frac{x^zP_r(\log x)}{1+x^2}\,dx.
\end{equation}
Now sum with multiplier $u^{r-1}$ over $r\geq1$. In a sufficiently
small neighborhood, for example $|z|+2|u|<1$, all the integrations and
sums involved are dominated by integrable powers of $x$ at their
endpoints. The two elementary generating identities are
\begin{align}
 \sum_{r\geq1}u^{r-1}\{\zeta(r,b)-\zeta(r)\}
   &=\psi(1-u)-\psi(b-u),\label{mld:gh:eq:zetasum}\\
 \sum_{r\geq0}P_r(L)u^r
   &=\frac{\pi u}{\sin\pi u}e^{2uL},
       \qquad P_0=1.\label{mld:gh:eq:Psum}
\end{align}
In \eqref{mld:gh:eq:zetasum}, the $r=1$ value is interpreted as above;
the formula follows by summing
$(n+b-u)^{-1}-(n+1-u)^{-1}$ over $n\geq0$.
Equation \eqref{mld:gh:eq:Psum} follows from
$\pi u/\sin\pi u=1+2\sum_{j\geq1}\eta(2j)u^{2j}$.

The sum of the left sides of \eqref{mld:gh:eq:integerreflection} is
$A(z,u)+A(-z,-u)$. Equation \eqref{mld:gh:eq:zetasum} gives the first
line of \eqref{mld:gh:eq:B}; equation \eqref{mld:gh:eq:Psum}, with $u$
replaced by $-u$, gives
\[
 \int_0^1\frac{x^z}{1+x^2}
   \left\{\frac{\pi x^{-2u}}{\sin\pi u}-\frac1u\right\}\,dx
 =\frac{\pi D(z-2u)}{\sin\pi u}-\frac{D(z)}u.
\]
This proves the claimed identity as an analytic germ.

The paired series in Lemma~\ref{mld:gh:lem:convergence} gives a meromorphic
continuation of $A$ to $\mathbb C^2$: near any candidate polar
hyperplane, multiply by its local linear factors; the same uniform
tail estimates apply. Its possible polar hyperplanes are
$z=3,5,\ldots$ and $u=1,2,\ldots$.
The right side of \eqref{mld:gh:eq:B} is also meromorphic on $\mathbb C^2$.
The identity theorem for meromorphic functions extends the equality
from the germ to the whole connected domain. The left side is
holomorphic on $\Omega$, so all apparent poles of the complete
right side there cancel. Finally, the sign in the Taylor projection
is $(-1)^{(p-1)+(r-1)}=(-1)^{p+r}$, proving the final assertion.
\end{proof}

\begin{corollary}[Every reflected shifted harmonic sum]
\label{mld:gh:cor:derivatives}
For positive integers $p,r$ and $(z,u)\in\Omega$,
\begin{align}
 &\sum_{n=1}^{\infty}(-1)^n
   \left\{\frac{h_n^{(r)}(u)}{(2n+1-z)^p}
       +(-1)^{p+r}\frac{h_n^{(r)}(-u)}{(2n+1+z)^p}\right\}
       \notag\\
 &\hspace{14mm}=
 \frac{\partial_z^{p-1}\partial_u^{r-1}B(z,u)}
      {(p-1)!(r-1)!}.\label{mld:gh:eq:allmixed}
\end{align}
The right side is a finite expression in polygamma functions and
elementary trigonometric functions. The equality holds for ordinary
conditionally convergent sums when $p=1$.
\end{corollary}
\begin{proof}
Differentiate \eqref{mld:gh:eq:reflection}. Lemma~\ref{mld:gh:lem:convergence}
justifies the operation, including $p=1$. The identities
$\partial_u^{r-1}h_n^{(1)}(u)=(r-1)!h_n^{(r)}(u)$ and
$\partial_z^{p-1}(2n+1-z)^{-1}=(p-1)!(2n+1-z)^{-p}$
give the first term. Applying the chain rule to $A(-z,-u)$ contributes
the factor $(-1)^{p+r}$.
\end{proof}

At rational parameter pairs in the domain, the right side uses only
polygamma values at rational arguments and elementary constants. The
digamma differences reduce by Gauss's formula to logarithms of algebraic
numbers and algebraic multiples of $\pi$. Positive polygamma orders
reduce to Hurwitz zeta values, and hence by residue-class filters to
depth-one polylogarithms at roots of unity. This gives an explicit
depth-one expression for the reflected combination; it does not assert
the same for either summand separately.

\subsection{A finite formula for the unshifted parity coefficients}

Write $\beta(s)=\sum_{n\ge0}(-1)^n/(2n+1)^s$ for the Dirichlet
beta function, initially on $\Re s>0$ and continued analytically.
For $j\geq0$, define
\[
 d_0(r)=
 \begin{cases}2\log2,&r=1,\\(2^r-2)\zeta(r),&r>1,\end{cases}
 \qquad
 d_j(r)=\frac{(r)_j(2^{r+j}-1)\zeta(r+j)}{2^j j!}
 \quad(j\geq1),
\]
and let
$c_{2\ell}=|E_{2\ell}|\pi^{2\ell}/(2^{2\ell}(2\ell)!)$, where
$\sec t=\sum_{\ell\geq0}|E_{2\ell}|t^{2\ell}/(2\ell)!$.
Put
\begin{align}
 C_{p,r}={}&2^r\binom{p+r-1}{r}\beta(p+r)\notag\\
 &+2\sum_{j=1}^{\lfloor r/2\rfloor}
      2^{r-2j}\eta(2j)
      \binom{p+r-2j-1}{r-2j}\beta(p+r-2j).
      \label{mld:gh:eq:Cpr}
\end{align}
The coefficient form of the reflection is
\begin{equation}\label{mld:gh:eq:finiteparity}
 \boxed{
 Q_{p,r}=\frac{(-1)^{r+1}}2
 \left\{C_{p,r}-\frac\pi2
   \sum_{\ell=0}^{\lfloor(p-1)/2\rfloor}
        c_{2\ell}d_{p-1-2\ell}(r)\right\},
 \qquad p+r\ \text{even}.}
\end{equation}
Indeed, extract $z^{p-1}$ in \eqref{mld:gh:eq:integerreflection}; the
coefficient on its left is $2Q_{p,r}$ in the indicated parity.
Taylor expansion of the Hurwitz zeta gives the $d_j$ terms.
Integration of each monomial in \eqref{mld:gh:eq:Pr} gives the beta
terms in \eqref{mld:gh:eq:Cpr}; the common sign is $(-1)^{r+1}$.
This is a finite coefficient algorithm for every $p,r$, with no
integer-relation search.

For example, writing $G=\beta(2)$,
\begin{align}
 Q_{2,2}&=\frac{7\pi}{4}\zeta(3)-6\beta(4)-\frac{\pi^2G}{12},
       \label{mld:gh:eq:Q22}\\
 Q_{1,3}&=4\beta(4)+\frac{\pi^2G}{6}
                     -\frac{3\pi}{2}\zeta(3),\\
 Q_{4,2}&=\frac{31\pi}{8}\zeta(5)+\frac{7\pi^3}{32}\zeta(3)
                     -20\beta(6)-\frac{\pi^2}{12}\beta(4),\\
 Q_{3,3}&=40\beta(6)+\frac{\pi^2}{2}\beta(4)
                     -\frac{93\pi}{8}\zeta(5)
                     -\frac{3\pi^3}{16}\zeta(3),\\
 Q_{4,4}&=-280\beta(8)-\frac{10\pi^2}{3}\beta(6)
                     -\frac{7\pi^4}{720}\beta(4)
                     +\frac{635\pi}{8}\zeta(7)
                     +\frac{31\pi^3}{16}\zeta(5).
\end{align}
Every term has total weight $p+r$. Only odd zeta values, even beta
values, $\pi$, and, in the $r=1$ case, $\log2$ occur.

\subsection{A rational-shift identity with an elementary answer}

\begin{corollary}\label{mld:gh:cor:quartershift}
The ordinary convergent nested sum satisfies
\begin{equation}\label{mld:gh:eq:quartershift}
 \boxed{
 \sum_{n=1}^{\infty}\frac{(-1)^n}{2n+1}
  \sum_{j=1}^n\left\{\frac1{j-\frac14}+\frac1{j+\frac14}\right\}
 =\pi\bigl(\log(1+\sqrt2)-1\bigr).}
\end{equation}
Equivalently, dividing by $32$, the inner summand is
$j/(16j^2-1)$ and the right side is
$\pi(\log(1+\sqrt2)-1)/32$.
\end{corollary}
\begin{proof}
Set $(z,u)=(0,1/4)$ in \eqref{mld:gh:eq:reflection}.
The digamma reflection formula gives
$\psi(3/4)-\psi(1/4)=\pi$, and $D(0)=\pi/4$.
The remaining resolvent is
\[
 D(-1/2)=2\int_0^1\frac{dt}{1+t^4}
 =\frac{\pi}{2\sqrt2}
       +\frac{\log(1+\sqrt2)}{\sqrt2}.
\]
The last integral follows by factoring
$1+t^4=(t^2+\sqrt2t+1)(t^2-\sqrt2t+1)$ and integrating the two
partial fractions. Substitution into \eqref{mld:gh:eq:B} gives
$-\pi^2/2+\pi\sqrt2D(-1/2)-\pi$, which is the stated answer.
\end{proof}

\subsection{Verification and the unresolved odd-weight direction}

The accompanying \texttt{derive\_parity.py} generates the finite symbolic
formula \eqref{mld:gh:eq:finiteparity}. It independently compares all $25$
ordered pairs with $p+r\leq10$ and $p+r$ even against real quadrature
of \eqref{mld:gh:eq:moment} at $60$ decimal working digits. The largest
absolute discrepancy is less than $3\times10^{-58}$.
Three separate $450$-term Euler evaluations of \eqref{mld:gh:eq:A}, including
a complex parameter pair and a point outside the initial small-parameter
proof neighborhood, agree with \eqref{mld:gh:eq:B} to more than $115$
decimal places. The JSON file records every tested pair and residual.
Three mixed-derivative checks, including $p=1,r=2$, and the quarter-shift
specialization are checked by the same independent series method; their
absolute discrepancies are below $3\times10^{-136}$.
These checks audit the signs and normalizations; the arbitrary-order
proof is the analytic argument above.

At $(z,u)=(0,0)$, equation \eqref{mld:gh:eq:allmixed} has coefficient
$1+(-1)^{p+r}$. It therefore gives no value for $Q_{p,r}$ when
$p+r$ is odd. In particular, $S_6=Q_{6,1}$ and $S_8=Q_{8,1}$ remain
outside the scalar parity projection. This is a precise limitation
of the reflection method, not a proof of arithmetic irreducibility.
The repository already exhibits a nonzero separating functional for
the restricted depth-two product presentation of the $S_6$ candidate;
the present identity does not supply the missing odd-weight relation.
